\begin{textsample}{POS Dim 2 – 2017\_c5 – Score -2.00 – 2017\_c5/enade\_2017\_admin\_1e32d435fe.txt}  \label{ex:f2_pos_392}
CAPÍTULO 5 DISTRIBUIÇÃO DOS CONCEITOS 5.1 PANORAMA NACIONAL DA DISTRIBUIÇÃO DOS CONCEITOS O Conceito Enade é um indicador de qualidade que avalia os cursos por intermédio dos desempenhos dos estudantes a partir dos resultados do Enade.
Os valores possíveis vão de 1 ( pior situação ) a 5 ( melhor situação ).
Os cursos com apenas um ou sem nenhum concluinte participante não obtêm o Conceito Enade, ficando Sem Conceito ( SC ).
A Tabela 5.1 apresenta a quantidade e distribuição de cursos de Arquitetura e Urbanismo participantes do Enade/2017, por faixa de conceito e Grande Região.
Enfatiza -se, mais uma vez, que a diferença entre os cursos tabulados neste capítulo e no capítulo 2 corresponde aos cursos Sem Conceito ( SC ), ou seja àqueles, em princípio, sem alunos concluintes que houvessem participado da prova ou àqueles que tivessem tão somente um aluno concluinte18.
Observando -se os dados da Tabela 5.1, nota -se que, dos 307 cursos participantes, 125 ( 40,7 \% ) classificaram -se com conceito 3, o valor modal.
Esse foi também o conceito modal nas regiões Sudeste ( 47,8 \% ) e Sul ( 46,1 \% ).
Na região Nordeste, os conceitos 2 e 3 foram igualmente atribuídos a 31,0 \% dos cursos.
Nas regiões Norte e Centro-Oeste, a moda foi o conceito 2, atribuído a 52,4 \% e 46,9 \% dos cursos da região, respectivamente.
O conceito 2 foi o segundo mais frequente em nível nacional ( 30,6 \%, correspondendo a 94 cursos ), e o conceito 4, o terceiro ( 18,9 \%, correspondendo a 58 cursos ).
Houve, ainda, 20 cursos ( 6,5 \% ) que receberam conceito 5, e sete cursos ( 2,3 \% ) que receberam conceito 1.
Dos 307 cursos de Arquitetura e Urbanismo, três ( 1,0 \% ) ficaram Sem Conceito ( SC ). 18 Esses cursos com somente um concluinte não foram considerados no capítulo 2.
A região Norte participou com 21 cursos ou 6,8 \% do total nacional.
Desses, 11 cursos receberam o conceito 2, o que equivale a 52,4 \% do total regional, além de ser o conceito modal na região, como já comentado.
A dois cursos ( 9,5 \% ) atribuiu -se conceito 1, a quatro cursos ( 19,0 \% ), conceito 3, e a outros quatro cursos ( 19,0 \% ), conceito 4.
Nenhum curso ficou Sem Conceito ( SC ) ou recebeu conceito 5 nessa região.
A região Nordeste participou com 42 cursos ou 13,7 \% do total nacional.
Nessa região, o conceito modal ficou dividido entre os conceitos 2 e 3, que apresentaram o mesmo número de cursos ( 13 ) e o mesmo percentual ( 31,0 \%, em termos regionais ).
O conceito 4 foi atribuído a dez cursos ( 23,8 \% ).
Ainda nessa região, seis cursos ( 14,3 \% ) foram avaliados com o conceito 5.
Nenhum curso ficou com conceito 1, assim como nenhum curso ficou Sem Conceito ( SC ).
Dos 136 cursos participantes da região Sudeste ( 44,3 \% do total nacional ), o conceito 3 foi o conceito modal, como já anteriormente assinalado, tendo sido obtido por 65 cursos ( 47,8 \% ).
O conceito 1 foi atribuído a dois cursos ( 1,5 \% ), e o conceito 2, a 39 cursos ( 28,7 \% ).
Receberam os conceitos 4 e 5, respectivamente, 22 cursos ( 16,2 \% ) e seis cursos ( 4,4 \% ).
Nessa região, dois cursos ( 1,5 \% ) ficaram Sem Conceito ( SC ).
Os 76 cursos da região Sul corresponderam a 24,8 \% do total nacional.
Como já indicado, houve predominância do conceito 3, atribuído a 46,1 \% dos cursos da região ( conceito modal ), ou seja, atribuído a 35 dos 76 cursos participantes na região Sul.
O conceito 1 foi atribuído a um curso ( 1,3 \% ), o conceito 2, a 16 cursos ( 21,1 \% ), o conceito 4, a 17 cursos ( 22,4 \% ), e o conceito 5, a seis cursos ( 7,9 \% ).
Um curso ( 1,3 \% ) ficou Sem Conceito ( SC ) na região Sul.
Já dos 32 cursos participantes na região Centro-Oeste ( 10,4 \% do total nacional ), 15 receberam conceito 2, o conceito modal.
A dois cursos ( 6,3 \% ) foi atribuído conceito 1, a oito cursos ( 25,0 \% ), o conceito 3, a cinco cursos ( 15,6 \% ), o conceito 4, e a dois cursos ( 6,3 \% ) o conceito 5.
Nessa região, nenhum curso ficou Sem Conceito ( SC ).
O Gráfico 5.1 apresenta a distribuição cumulativa do Conceito Enade, segundo Grande Região.
Num gráfico de distribuição cumulativa, em cada valor de abscissa ( eixo x ), o valor apresentado no eixo das ordenadas ( eixo y ) é o valor acumulado da variável.
Por exemplo, a região Norte ( linha azul ) apresenta um pouco mais de 80 \% dos cursos com conceito 3 ou abaixo deste valor.
A região Nordeste ( linha verde ), por outro lado, apresenta 61,9 \% dos cursos na mesma situação : com conceito 3 ou abaixo.
Na região Norte ( linha azul ), os cursos têm conceitos 4 ou abaixo disso, já que a poligonal cumulativa alcança 100 \% para o conceito 4.
Em linhas gerais, num gráfico de distribuição cumulativa, poligonais mais à esquerda representam Grandes Regiões com pior distribuição, e poligonais mais à direita, Grandes Regiões com melhores conceitos.
A região Norte ( linha azul ) apresenta os cursos com conceitos mais baixos, ao passo que as regiões Nordeste ( linha verde ) e Sul ( linha roxa ) apresentam as distribuições com valores maiores.
Cursos em IES situadas nas regiões Sudeste ( linha vermelha ) e Centro-Oeste ( linha preta ) apresentam uma situação intermediária. 5.2 CONCEITOS POR CATEGORIA ADMINISTRATIVA, POR MODALIDADE DE ENSINO E GRANDE REGIÃO A Tabela 5.2 apresenta a distribuição dos cursos de Arquitetura e Urbanismo participantes do Enade/2017, por Categoria Administrativa e Modalidade de Ensino, de acordo com os conceitos por eles alcançados, segundo as Grandes Regiões brasileiras.
Dos 307 cursos participantes, 57 ( 18,6 \% ) eram ministrados em Instituições Públicas, e 250 ( 81,4 \% ), em Instituições Privadas.
Quanto à Modalidade de Ensino todos os cursos eram Presenciais, e nenhum era a Distância.
De acordo com as informações da Tabela 5.2, em termos nacionais, entre Instituições Públicas e Privadas, observa -se que, dos 20 cursos avaliados com conceito 5, 19 eram oferecidos em IES Públicas, e apenas um, em IES Privadas.
Dos 57 cursos participantes de IES Públicas, o conceito 4 foi o valor modal, atribuído a 27 cursos ( 47,4 \% ).
Entre os demais cursos de Instituições Públicas participantes, dois cursos ( 3,5 \% ) foram avaliados com conceito 1, outros dois cursos ( 3,5 \% ) obtiveram conceito 2, sete cursos ( 12,3 \% ) receberam conceito 3 e, como já comentado, 19 cursos ( 33,3 \% ) receberam conceito 5.
Nesta categoria, nenhum curso ficou Sem Conceito ( SC ).
Na rede privada, o conceito modal foi 3, atribuído a 118 cursos, o correspondente a 47,2 \% dos 250 cursos da categoria.
Entre os demais cursos participantes, cinco ( 2,0 \% ) receberam conceito 1, e 92 ( 36,8 \% ), conceito 2.
O conceito 4 foi atribuído a 31 cursos ( 12,4 \% ), e o conceito 5, a um curso ( 0,4 \% ).
Nesta Categoria Administrativa, três cursos ( 1,2 \% ) ficaram Sem Conceito ( SC ).
Dos 307 cursos participantes foram na modalidade Presencial, e o conceito 3, o valor modal, foi atribuído a 125 cursos ( 40,7 \% ).
Entre os demais cursos nesta Modalidade de Ensino, 94 cursos ( 30,6 \% ) receberam conceito 2, 58 cursos ( 18,9 \% ) obtiveram conceito 4, 20 cursos ( 6,5 \% ) receberam conceito 5, e sete cursos ( 2,3 \% ) conceito 1.
Nessa modalidade, três cursos ( 1,0 \% ) ficaram Sem Conceito ( SC ).
Como já mencionado, não houve cursos na modalidade Educação a Distância.
Na análise por região, observa -se que, na região Norte, as Instituições Privadas participaram com 16 cursos ( 76,2 \% do total regional ).
Nenhum ficou Sem Conceito ( SC ), e nenhum recebeu os conceitos 4 e 5.
O conceito modal para as Instituições Privadas na região foi 2, atribuído a 11 cursos, correspondendo a 68,8 \% dessa categoria na região.
Dois cursos ( 12,5 \% ) foram avaliados com conceito 1, três cursos ( 18,8 \% ) com conceito 3.
As Instituições Públicas participaram com cinco cursos na região Norte ( 23,8 \% do total regional ), dos quais quatro ( 80,0 \% ) obtiveram conceito 4, o valor modal.
O conceito 3 foi atribuído ao outro curso ( 20,0 \% ).
Nenhum curso ficou Sem Conceito ( SC ) ou com os conceitos 1, 2 e 5 nessa combinação de categoria e região.
Todos os cursos oferecidos na região Norte eram Presenciais, não havendo cursos na modalidade a Distância.
Na região Nordeste, a rede privada concentrou 30 dos 42 cursos participantes, o equivalente a 71,4 \% do total da região.
Dentre os cursos oferecidos por IES Privadas no Nordeste, 13 cursos ( 43,3 \% ) foram avaliados com conceito 3, o conceito modal.
Nenhum curso obteve conceito 1, 12 cursos ( 40,0 \% ) receberam conceito 2, quatro cursos ( 13,3 \% ) ficaram com conceito 4, e um curso ( 3,3 \% ) alcançou o conceito máximo, 5.
Nessa combinação de categoria e região, nenhum curso ficou Sem Conceito ( SC ).
As Instituições Públicas da região Nordeste participaram com 12 cursos ( 28,6 \% do total da região ).
Desses, seis ( 50,0 \% ) obtiveram conceito 4, o valor modal, seguidos por outros cinco cursos ( 41,7 \% ), que receberam conceito 5.
O conceito 2 foi atribuído a um curso ( 8,3 \% ).
Nenhum curso ficou Sem Conceito ( SC ) ou com conceitos 1 ou 3 nessa região.
Todos os cursos na região Nordeste foram Presenciais : 42.
Para esses, a moda foi contígua entre os conceitos 2 e 3, 13 cursos cada ( 31,0 \% ).
Dez cursos ( 23,8 \% ) receberam conceito 4 e seis cursos ( 14,3 \% ) receberam conceito 5.
Nenhum curso Presencial ficou com conceito 1 ou Sem Conceito ( SC ).
Na região Sudeste, a proporção de cursos da rede privada, 87,5 \%, foi mais elevada do que a registrada nas demais regiões brasileiras, correspondendo a 119 dos 136 cursos participantes.
Nessa categoria e região, o conceito modal foi 3, atribuído a 63 cursos ( 52,9 \% ).
Dois cursos ficaram Sem Conceito ( SC ).
Os demais foram avaliados com conceito 1 ( um curso ), conceito 2 ( 39 cursos ), conceito 4 ( 14 cursos ) e nenhum curso recebeu o conceito 5.
Entre os 17 cursos oferecidos por Instituições Públicas na região Sudeste, o conceito modal foi 4, atribuído a oito cursos ( 47,1 \% da região ).
Os demais cursos receberam os conceitos 1 ( um curso ), 3 ( dois cursos ) e 5 ( seis cursos ).
Na região Sudeste, nesta categoria, nenhum curso recebeu conceito 2 ou ficou Sem Conceito ( SC ).
Também na região Sudeste, todos os cursos foram Presenciais : 136.
Para esses, a moda foi o conceito 3 com 65 cursos ( 47,8 \% ).
Os demais cursos Presenciais da região foram alocados aos conceitos 1 ( dois cursos ), 2 ( 39 cursos ), 4 ( 22 cursos ) e 5 ( seis cursos ).
Somente dois cursos ficaram Sem Conceito ( SC ).
As Instituições Privadas concentraram 60 dos 76 cursos participantes da região Sul, 78,9 \% do total regional.
Desses, 32 ( 53,3 \% ) obtiveram conceito 3, o conceito modal.
Nessa combinação de Categoria Administrativa e Grande Região, os conceitos 2 e 4 foram atribuídos a 15 e a 12 cursos, respectivamente.
Nenhum curso recebeu conceito 5, assim como nenhum curso recebeu conceito 1.
Apenas um curso ficou Sem Conceito ( SC ).
As Instituições Públicas na região Sul participaram com 16 cursos, 21,1 \% dos cursos da região.
O conceito modal foi 5, atribuído a seis cursos ( 37,5 \% ).
Aos demais cursos foram atribuídos os conceitos 4 ( cinco cursos ), 3 ( três cursos ), 2 ( um curso ) e 1 ( um curso ).
Nenhum curso da região Sul oferecido por IES Pública ficou Sem Conceito ( SC ).
Também na região Sul, todos os cursos foram Presenciais.
Para esses cursos, a moda foi o conceito 3 com 35 cursos ( 46,1 \% ).
Os demais cursos Presenciais da região foram alocados aos conceitos 1 ( um curso ), 2 ( 16 cursos ), 4 ( 17 cursos ) e 5 ( seis cursos ).
Somente um curso ficou Sem Conceito ( SC ).
Na região Centro-Oeste, 25 dos 32 cursos participantes eram de Instituições Privadas ( 78,1 \% em termos regionais ).
Desses, 15 receberam conceito 2, o conceito modal.
Entre os demais, dois cursos receberam conceito 1, sete receberam conceito 3 e um recebeu conceito 4.
Nenhum curso recebeu conceito 5 ou ficou Sem Conceito ( SC ).
Quanto aos sete cursos oferecidos por Instituições Públicas na região Centro-Oeste ( 21,9 \% do total regional ), quatro cursos apresentaram conceito 4, conceito modal, dois ficaram com conceito 5, e apenas um curso ficou com conceito 3.
Nenhum curso ficou Sem Conceito ( SC ) ou com conceitos 1 e 2.
Nessa região, todos os cursos foram Presenciais.
Para esses cursos Presenciais, a moda foi o conceito 2 com 15 cursos ( 46,9 \% ).
Os demais foram alocados aos conceitos 1 ( dois cursos ), 3 ( oito cursos ), 4 ( cinco cursos ) e 5 ( dois cursos ).
Nenhum curso Presencial ficou Sem Conceito ( SC ) nessa região.
O Gráfico 5.2 apresenta a distribuição do Conceito Enade, segundo a Categoria Administrativa da IES.
Os cursos em IES Públicas ( linha azul ) apresentam uma poligonal mais à direita do que os em IES Privadas ( linha verde ) e, consequentemente, uma distribuição de Conceitos Enade com valores maiores. 5.3 CONCEITOS POR ORGANIZAÇÃO ACADÊMICA E GRANDE REGIÃO Na Tabela 5.3, encontra -se a distribuição dos conceitos atribuídos aos cursos participantes do Enade/2017 na Área de Arquitetura e Urbanismo, por Organização Acadêmica, segundo as Grandes Regiões brasileiras.
Dos 307 cursos de Arquitetura e Urbanismo participantes, 161 eram oferecidos em Universidades, 87 em Centros Universitários, 57 em Faculdades e dois em CEFET/IFET.
Essa distribuição corresponde a, respectivamente, 52,4 \%, 28,3 \%, 18,6 \% e 0,7 \% dos cursos.
De acordo com os dados apresentados, de todos os 20 cursos avaliados com conceito 5, 19 eram vinculados a Universidades.
Os cursos vinculados a esse tipo de Organização Acadêmica tiveram o conceito 3 como conceito modal, atribuído a 59 cursos ( 36,6 \% ).
Os demais cursos vinculados a Universidades receberam os conceitos 1 ( três cursos ), 2 ( 37 cursos ), 4 ( 43 cursos ) e 5 ( 19 cursos, como já mencionado ).
Nenhum curso ficou Sem Conceito ( SC ).
Entre os cursos em Centros Universitários, o conceito modal foi também 3, atribuído a 46 cursos ( 52,9 \% ).
Vinculados a esse tipo de Organização Acadêmica, dois cursos ficaram Sem Conceito ( SC ).
Os outros cursos receberam os conceitos 1 ( dois cursos ), 2 ( 29 cursos ) e 4 ( 8 cursos ).
Nenhum curso recebeu conceito 5.
Dos 57 cursos mantidos por Faculdades, 28 ( 49,1 \% ) receberam conceito 2, o conceito modal.
O segundo conceito mais frequentemente atribuído aos cursos oferecidos por essa categoria de IES foi o conceito 3, recebido por 20 cursos.
Na sequência, dois cursos obtiveram conceito 1, cinco cursos obtiveram conceito 4, e um curso, conceito 5.
Apenas um curso ficou Sem Conceito ( SC ).
Dos dois cursos mantidos por CEFET/IFET, os dois receberam conceito 4.
Consequentemente, nenhum curso ficou Sem Conceito ( SC ) ou recebeu os conceitos 1, 2, 3 ou 5.
Considerando -se separadamente as regiões brasileiras, verifica -se que, na região Norte, as Universidades participaram com sete dos 21 cursos participantes.
Quatro cursos receberam conceito 4, conceito modal.
Os demais cursos receberam os conceitos 2 ( um curso ) e 3 ( dois cursos ).
Nenhum curso ficou Sem Conceito ( SC ) ou recebeu conceitos 1 e 5 nessa combinação de categoria e região.
Os Centros Universitários da região Norte foram representados por quatro cursos, aos quais foram atribuídos os conceitos 2 ( valor modal, recebido por dois cursos ), 1 e 3 ( recebidos por um curso cada ).
Nenhum curso ficou Sem Conceito ( SC ) ou recebeu os conceitos 4 e 5 nessa combinação de categoria e região.
As Faculdades participaram com dez cursos na região Norte, que receberam os conceitos 2 ( oito cursos, conceito modal ), 1 e 3 ( um curso cada ).
Nenhuma Faculdade da região ficou Sem Conceito ( SC ) ou obteve conceitos 4 ou 5.
Os CEFET/IFET não tiveram cursos participantes.
Na região Nordeste, as Universidades participaram com 19 dos 42 cursos da Área de Arquitetura e Urbanismo oferecidos na região, dos quais nenhum ficou Sem Conceito ( SC ).
Sete cursos obtiveram conceito 4, que foi o modal.
Quatro cursos obtiveram o conceito 2, três cursos obtiveram o conceito 3, e cinco cursos, o conceito 5.
Nenhum curso oferecido por Universidades no Nordeste recebeu conceito 1.
Os Centros Universitários contaram com 12 cursos participantes na região Nordeste.
Seis cursos obtiveram 3, o conceito modal.
Cinco cursos obtiveram conceito 2, e apenas um curso obteve conceito 4.
Nenhum curso oferecido por Centros Universitários no Nordeste ficou Sem Conceito ( SC ) ou recebeu conceitos 1 ou 5.
As Faculdades foram representadas por 11 cursos na região Nordeste.
Nota -se que tais cursos apresentaram distribuição bimodal, uma vez que as maiores frequências se concentraram em dois conceitos, os conceitos 2 e 3, atribuídos, cada qual, a quatro cursos.
O conceito 4 foi atribuído a dois cursos e o conceito 5 foi atribuído a um curso.
Na região Nordeste, nenhum curso oferecido por Faculdades recebeu conceito 1, tampouco ficou Sem Conceito ( SC ).
Os CEFET/IFET não tiveram representantes na região.
Na região Sudeste, as Universidades concentraram 73 dos 136 cursos de Arquitetura e Urbanismo da região.
Entre os cursos em Universidades, nessa região, o conceito modal foi 3, atribuído a 33 cursos.
Os demais cursos receberam os conceitos 1 ( dois cursos ), 2 ( 18 cursos ), 4 ( 14 cursos ) e 5 ( seis cursos ).
Nenhum curso ficou Sem Conceito ( SC ).
Os Centros Universitários participaram com 42 cursos na região Sudeste, dos quais 24 obtiveram o conceito modal 3.
O conceito 2 foi atribuído a 13 cursos, e o conceito 4, a três cursos.
Nenhum curso obteve conceitos 1 ou 5.
Dois cursos ficaram Sem Conceito ( SC ) nessa combinação de Organização Acadêmica da região.
Já as Faculdades foram representadas por 19 cursos na região Sudeste, dos quais nenhum ficou Sem Conceito ( SC ) ou recebeu conceito 1 ou 5.
Os conceitos 2 e 3 foram os modais, com oito cursos cada.
O conceito 4 foi atribuído a três cursos.
Os CEFET/IFET participaram com dois cursos, os quais receberam conceito 4.
Sendo assim, nenhum curso ficou Sem Conceito ( SC ), ou recebeu os demais conceitos.
Dos 76 cursos da região Sul, 48 eram de Universidades, para os quais o conceito modal foi 3, atribuído a 18 cursos.
Nesse tipo de organização, nenhum curso ficou Sem Conceito ( SC ) ou recebeu conceito 1.
Os demais cursos receberam os conceitos 2 ( dez cursos ), 4 ( 14 cursos ) e 5 ( seis cursos ).
Dos 17 cursos participantes de Centros Universitários da região Sul, nenhum ficou Sem Conceito ( SC ) ou recebeu conceito 5.
O conceito modal, 3, foi atribuído a 11 cursos.
Os conceitos 1, 2 e 4 foram atribuídos a um, a dois e a três cursos, respectivamente.
Quanto aos 11 cursos vinculados a Faculdades na região Sul, seis receberam o conceito modal 3.
O conceito 2 foi atribuído a quatro cursos.
Um curso ficou Sem Conceito ( SC ) e nenhum recebeu conceito 1, 4 ou 5.
Nessa região não foram avaliados cursos em CEFET/IFET.
Na região Centro-Oeste, 14 dos 32 cursos eram de Universidades.
Nesse tipo de organização, nenhum curso ficou Sem Conceito ( SC ).
Os conceitos 2 e 4 foram atribuídos a quatro cursos cada, sendo esses os conceitos modais.
O conceito 3 foi atribuído a três cursos, o conceito 5 a dois cursos, e o conceito 1, a apenas um curso.
Os Centros Universitários da região Centro-Oeste contaram com 12 cursos, que obtiveram os conceitos 2 ( setes cursos, conceito modal ), 3 ( quatro cursos ) e 4 ( um curso ).
Nenhum curso obteve os conceitos 1 ou 5, assim como também nenhum curso ficou Sem Conceito ( SC ).
Já as Faculdades foram representadas por seis cursos na região Centro-Oeste, dos quais nenhum ficou Sem Conceito ( SC ) ou recebeu conceito 4 ou 5.
O conceito 2 foi o conceito modal, com quatro cursos e os conceitos 1 e 3 receberam apenas um curso cada.
Nessa região não foram avaliados cursos em CEFET/IFET.
O Gráfico 5.3 apresenta a distribuição do Conceito Enade dos cursos de Bacharelado, segundo a Organização Acadêmica da IES.
Os cursos em CEFET/IFET ( linha roxa ) apresentam uma poligonal mais à direita do que as demais, denotando uma melhor distribuição de Conceitos Enade.
No outro extremo, temos os cursos oferecidos em Faculdades ( linha vermelha ) com a poligonal mais à esquerda e, consequentemente, uma distribuição de Conceitos Enade com menores valores.
IES de Centros Universitários ( linha verde ) e em Universidades ( linha azul ) apresentam uma situação intermediária.

% matched lemmas: 
\end{textsample}
